% frontmatter/english/como-usar.tex -- "How to use this book": las
% instrucciones para el adulto de frontmatter/como-usar.tex, en inglés.
\section*{How to use this book}

This book is for a child of seven or eight who already reads on their
own and can add and take away up to 20. It does not teach grown-up
sums, but the ideas of everyday economics: that things have value, that
money lets us swap one thing for another, that we cannot have
everything and so we have to choose, that saving means keeping
something today for tomorrow, that somebody makes the bread we buy, and
that we all pay together for the pavements and the parks. It belongs to
the same family as \emph{First Words} and \emph{First Numbers}, with the
same characters (Lucía, her little brother Dani, their dog Toby,
Grandma Rosa, and Lucía's friends Sofía and Hugo). And it is, page by
page, the English edition of \emph{Aprendo economía}: the same story,
the same pictures and the same activity every day, so a child can use
either book, or both.

There is a page for every weekday of the year, Monday to Friday,
holidays included. Ten minutes a day is plenty. Every page works the
same way:

\begin{itemize}[leftmargin=6mm]
\item \textbf{Date}, \textbf{Done} and \textbf{Notes}: for the grown-up.
  Ticking each day whether it was done alone, with help or with
  difficulty takes a second, and by the end of the year it is a record
  of everything that has been learnt.
\item \textbf{Today}, in the green box: the day's story, two or three
  sentences for the child to read. On Mondays, underneath, the
  \textbf{word of the week} and what it means; on Friday it comes back,
  and at the end of the book they are all together in \emph{My
  economics dictionary}.
\item An \textbf{activity}. The instructions (in small grey print) can
  be read out by a grown-up, if needed.
\end{itemize}

Each activity is explained, one by one, on the next page.

\subsection*{How to help}

{\small
\begin{itemize}[leftmargin=6mm, itemsep=1pt, topsep=2pt]
\item Talk about the day's story: what would you have done? In
  economics there is hardly ever just one answer, and what children
  learn is to think about what they gain and what they give up.
\item Sums with money make more sense with real coins: count the ones
  you have at home, play shops, and let them pay and check the change
  when you go shopping.
\item Getting it wrong with a little money is the best way to learn to
  use it well, for when there is more.
\item The money in this book is euros, as in Spain, where the family
  lives, and always whole euros. The euro sign goes before the amount,
  as with pounds: €5 is five euros.
\item At the end of each part there is a medal, and at the end of the
  book, a diploma and the \textbf{answers} for the pages that have one.
\end{itemize}}
\newpage
